\scanpage{20}
% archive foliation: 19
2) Pour que $L_*$ soit \deleted{projectif}, \inserted{flasque (projectif)},
il f. et s.
qu'il n'y ait que des facteurs
$(\Lambda^i\Phi_*)^{\mu_i}$
\inserted{(resp. et $(\Lambda^0\Psi_*)^{\nu_0}=k^{\nu_0}$)}.
Pour que $L_*\otimes_k K$ (comme
$K$-module $\frac12$ simplicial, où $K$ est le
corps des fractions) soit flasque
il f. et suffit qu'il n'y ait que des
$(\Lambda^i\Phi_*)^{\mu_i}$, $\Phi_*(i,n)^{\beta_{i\ell}}$
\inserted{$(i\geq1)$}, à l'exclusion
des $(\Lambda^j\Psi_*)^{\nu_j}$ $(j\geq0)$.

\deleted{Pour}
3) Pour déterminer d'ailleurs toute la
catégorie additive des $k$-modules
$\frac12$ simpliciaux str. $\partial$-ps coh. (resp.
str. $\partial$-parfaits), il suffit en principe
de déterminer les hom. modules
entre les objets indécomposables
\[
\Lambda^i\Phi,\quad\uline{\Phi(i,n)},\quad\Lambda^j\Psi_*
\qquad(i\geq1,j\geq0)
\]
\inserted{\unsure{un pour chaque idéal}\ \unsure{chaque diviseur}}
qu'on peut réduire à deux types si
on admet aussi le cas $n=1$
\[
\Phi(i,n)\qquad\Lambda^j\Psi_*
\qquad(i\geq1,j\geq0)
\]
On trouve, en repassant aux
complexes de chaînes associés
